\section{Each Element in Its Setting}
\label{sec:each-element}

Seven of the ten elements are Scripture and three are the Church's own orations.
The Introit joins a prayer from the Greek chapters of Esther to the first verse
of Ps~118; the Gradual, the Alleluia and the Communion take verses from three
psalms, the Communion returning to Ps~118; the lessons are St Paul's letter to
the Ephesians and St Matthew's Gospel; and the Offertory tells the opening of
the Book of Job in words of its own. Each element is set out below in its own
context, with what its words say and what the liturgical commentators who
expound it say of it, before the three readings of the whole.

\subsection{Introit: Mardochai's prayer, and the blessed way}

The Book of Esther tells how a Jewish girl of the exile became queen of
Assuerus, king of Persia, at Susa, and how her cousin and guardian Mardochai
would not bow down to Aman, the king's favourite. Aman obtained a royal decree
that every Jew in the empire should be put to death on one day. The Hebrew book
tells of Mardochai's mourning in sackcloth and ashes and Esther's fast before
she went in to the king unbidden; it gives no prayer. The prayers of Mardochai
and of Esther stand in the chapters that the Vulgate prints after the end of the
Hebrew book (Esth 10:4--16:24), which St Jerome translated from the Greek and
set apart; the \work{Catholic Encyclopedia}'s article ``Esther'' reports them as
most probably translations from a Hebrew or Chaldaic original of which no trace
remains in the Hebrew Bible. They are part of the Book of Esther in the Church's
canon. The first of these chapters opens with a colophon that dates the bringing
of ``this epistle of Phurim'' to the fourth year of the reign of a Ptolemy and
Cleopatra (11:1).

``Mardochai besought the Lord, remembering all his works'' (13:8). His prayer
begins with the confession that the Introit takes: ``O Lord, Lord, almighty
king, for all things are in thy power, and there is none that can resist thy
will, if thou determine to save Israel'' (13:9). He confesses God the maker of
heaven and earth and Lord of all (13:10--11). Then, in lines the Mass does not
sing, he tells God that it was not pride that kept him from worshipping Aman,
for he would have kissed the steps of his feet for Israel's salvation, ``But I
feared lest I should transfer the honour of my God to a man'' (13:12--14). He
asks mercy on a people whose enemies ``resolve to destroy us'' (13:15), reminds
God of ``thy portion, which thou hast redeemed for thyself out of Egypt''
(13:16), and asks that God turn their mourning into joy, ``that we may live and
praise thy name, O Lord, and shut not the mouths of them that sing to thee''
(13:17). ``And all Israel with like mind and supplication cried to the Lord,
because they saw certain death hanging over their heads'' (13:18). The Missal's
antiphon keeps the confession and drops the condition ``if thou determine to
save Israel''. What Mardochai said of God's power to save one threatened people
the Introit says of God's will without condition: in it all things are placed,
and no one resists it.

No Father of the first centuries wrote a commentary on Esther. St Rabanus Maurus
expounded the book in the ninth century, from the Hebrew chapters only, passing
over the Greek additions, which included this prayer. Yet when he comes to the
king's edict hung up in Susa he reads it as the decree of the heavenly king,
whose rule reaches the whole world, and he describes that king in the words of
Mardochai's prayer, citing Esther~13: in his power all things are placed, and
none can resist his will. The pride of the world, he concludes, is overcome by
the power of God, as Christ said, ``I have overcome the world'' (Jn 16:33). Aman, whose name
Rabanus interprets as iniquity, stands for the powerful of this world who
persecute those who will not give them divine honour, and prefigures the enemies
of the Church (\work{Expositio in librum Esther}, preface, VI, XI and XIII). St
Thomas Aquinas takes the same verse as his authority that nothing can happen
outside the order of the divine government (\work{Summa theologiae} I q.~103
a.~7, \latin{sed contra}), an argument the first reading follows.

The liturgical commentators hear the antiphon in more than one voice. For
William Durandus the Church speaks with the voice of Mardochai, who prayed in
sackcloth and ashes, the arms of penance; and because God is Lord of all, he
frees those who serve him from the hand of the devil and helps them in their war
against him, as he freed Mardochai from the hand of Aman (\work{Rationale}
VI.138). Rupert of Deutz hears the voice of the soldier of the virtues, who in
his trial falls to the earth with Job and says, ``The Lord gave, and the Lord
hath taken away'' (\work{De divinis officiis} XII.21). Honorius Augustodunensis
reads the Sunday's chants from the history of Esther, through which God's people
is delivered from the danger hanging over it (\work{Gemma animae} IV.90). Lucien
Fromage, continuing Dom Gu\'eranger's \work{The Liturgical Year}, says that the
holy men who live in a world reduced, like Job on the dunghill, to the extremity
of wretchedness ``begin by making their own the sublime prayer made, by
Mardochai, for his people, who were doomed to extermination''. Bl.\ Ildefonso
Schuster draws from the same prayer what he calls the secret of the Church's
history: the evil done by the malice of creatures with God's permission enters
the plan of Providence, which reaches its highest ends through its bitterest
enemies, ``Like the devil, who, despitefully winnowing the grain, renders a
helpful service to the great Householder'' (\work{The Sacramentary} III,
pp.~178--179).

The verse is the first of Ps~118, the longest of the psalms, whose
twenty-two strophes of eight verses follow the Hebrew alphabet and praise God's
law. It opens with a beatitude, ``Blessed are the undefiled in the way, who walk
in the law of the Lord.'' St Augustine begins his sermons on the psalm from the
desire for blessedness that all men share: ``if thou wouldest be blessed, be
undefiled''; and he asks in what way, ``save in the law of the Lord?'' Yet he
will not let the undefiled forget that they are debtors: the saints who walk in
the Lord's ways still say, ``If we say that we have no sin, we deceive
ourselves'' (1~Jn 1:8; \work{Enarrationes in Psalmos} 118, sermon~1.1--3).
Cassiodorus notes that \latin{immaculati} is plural and said of the faithful,
not of Christ, who alone is spotless by nature, and that the saints become
undefiled when they receive the gifts of pardon, \latin{quando indulgentiae
munera consequuntur}; the way is the Saviour himself, who said, ``I am the way''
(\work{Expositio Psalmorum}, Ps~118:1).

\subsection{Collect: the family kept by continual mercy}

The Collect is the Church's own composition. It calls the Church God's family,
\latin{familiam tuam}, and asks two things of his \latin{continua pietas}: that
under his protection the family be free from every adversity, \latin{a cunctis
adversitatibus}, and that it be devoted to his name in good works. Its
\latin{te protegente} and its adversities stand beside the Gradual's refuge and
the Epistle's adversaries; its \latin{pietas} returns in the Secret's
\latin{potenti pietate}. Honorius reads it as the prayer of God's people in the
history of Esther, asking that God guard his family (\work{Gemma animae}
IV.90). Fromage finds in it the Church's preference for peace, ``because, that
furnishes her with undisturbed freedom for paying to her God the united homage of
religion and good works'', and joins it to the close of Mardochai's prayer,
which asked that the people might live to praise God's name. Schuster hears two
requests: protection against the wiles of the devil, who raises up obstacles
material, moral and spiritual to turn men from a righteous life, and the spirit
of \latin{pietas}, the gift of the Holy Spirit by which the mind is drawn to the
service of God and to good works (\work{The Sacramentary} III, p.~179).

\subsection{Epistle: the armour of God}

The lesson is the close of the letter to the Ephesians. Paul writes as a
prisoner, ``the prisoner of Jesus Christ'' (3:1), ``a prisoner in the Lord''
(4:1), and calls himself in
the verses that follow the lesson ``an ambassador in a chain'' (6:20). The
letter's last instruction is a call to arms. ``Be strengthened in the Lord''
places the strength in God; the armour is ``the armour of God''; the enemy is the
devil, whose ``deceits'', \latin{insidias}, are what the Christian must stand
against. Then comes the sentence on which the lesson turns: ``our wrestling is
not against flesh and blood; but against principalities and powers, against the
rulers of the world of this darkness, against the spirits of wickedness in the
high places.'' The Christian is to resist ``in the evil day'' and, having done
all, to stand. The armour is named piece by piece: the girdle of truth, the
breastplate of justice, the shoes of the gospel of peace, the shield of faith
that quenches ``all the fiery darts of the most wicked one'', the helmet of
salvation, and the sword of the Spirit, ``which is the word of God''. In the
letter the call goes on to prayer ``at all times in the spirit'' (6:18), which
the Mass does not read.

St John Chrysostom notes that Paul says not ``fightings'' but ``wiles'': ``this
enemy is at war with us, not simply, nor openly, but by `wiles'''; the devil
``never proposes to us sins in their proper colors''. ``By `evil day' he means
the present life''; to stand is to stand ``after we have slain'', for many who
won have fallen again (\work{Homilies on Ephesians} 22). St Jerome reads the
whole armour as Christ himself: by the armour of God, he says, we are to
understand the Saviour, so that ``Put on the armour of God'' says the same as
``Put on the Lord Jesus Christ'' (Rom 13:14): \latin{quasi dixerit, Induite
Dominum Jesum Christum} (\work{Commentarii in Epistolam ad Ephesios} III, PL~26,
col.~543). He adds that the wars of Israel against the nations, David against
Goliath, Josue, the Judges and the Kings, were written as a figure of this
spiritual war. St Thomas Aquinas makes the armour the gifts and virtues, the
girdle temperance, the breastplate justice and the shield faith; and to the
objection that men plainly do attack us he answers that flesh and blood attack
us not principally but as moved by a higher agent, the devil, \latin{a superiore
movente} (\work{Super Epistolam ad Ephesios} 6, lect.~3--4).

Durandus says that even the well armed cannot perfectly resist unless God
helps them, and finds there the reason for the Gradual's refuge (\work{Rationale}
VI.138); Rupert hears the Epistle comfort the soldier whom the Introit showed
fallen with Job (XII.21). Fromage, like Jerome, makes the armour first Christ's:
the prophets describe the Leader himself girded, helmeted, in breastplate and
shield and with the sword of the spirit, so that ``We should almost think we were
here having a list of our own arms; well, yes, but they are his, first'' (Is
11:5; 59:17; Wis 5:19--20; Apoc 2:16). The evil days are many in
every life, he says, but ``there is one evil day, evil beyond all the others: it
is the last day, the day of judgment'' (\work{The Liturgical Year} XI,
pp.~459--460). Schuster calls faith the most effectual defence against the
devil, and says that God allows some of the saints to sustain this combat ``in
order that their faith may be still further purified and strengthened'', ``This
is precisely the history of Job'' (\work{The Sacramentary} III, p.~179).

\subsection{Gradual: the eternal God our refuge}

The Gradual sings the first two verses of Ps~89, which the Vulgate entitles
\latin{Oratio Moysi, hominis Dei}, ``A prayer of Moses the man of God''. The
title is not sung. The psalm goes on, in verses the Mass leaves aside, to set
the shortness and misery of human life against God's eternity: a thousand
years in his sight are as yesterday (v.~4), and man is grass that flourishes in
the morning and withers by evening. The two verses the Gradual sings speak only
of God: he has been our refuge \latin{a generatione et progenie}, and before the
mountains were made, or the earth and the world formed, \latin{a saeculo et usque
in saeculum tu es, Deus}.

St Augustine, reading the title, calls Moses ``the Minister of the Old, and the
Prophet of the New Testament'', and takes ``from one generation to another'' as
either every generation or the two, the old and the new. He does not believe the
psalm ``was entirely the composition of that Moses'', since it lacks the turns
of phrase of Moses' songs, and holds that the great servant's name is used ``for
the sake of some intimation''. On the second verse he observes that the psalmist
does not say ``thou wast'' and ``thou shalt be'' but puts the verb in the
present, ``intimating that the substance of God is altogether immutable'', and
concludes: ``Behold then the eternity that is our refuge, that we may fly thither
from the mutability of time, there to remain for evermore'' (\work{Enarrationes}
89.1--3). St Robert Bellarmine, in his \work{Explanation of the Psalms}, notes
that the Hebrew word for refuge ``signifies a well fortified house, placed on an
eminence'', and that those who dwell in God dwell ``in a city fortified by faith,
hope, and charity''; he too marks the present tense, ``Thou art God'', ``in order
to show the true eternity of God, in which there is no past or future''. On the
title's ascription to Moses, the Doctor admits in his general prefaces that
those named in the titles of the psalms may have written them, but holds the
more probable opinion to be that David wrote them all, as he says Augustine
maintains. He and Augustine differ only in that Bellarmine leaves Moses'
authorship open as a tenable opinion; the title itself names Moses, and the
Scriptural Date and Location appendix sets out the witnesses.

Cassiodorus reads the first verse through the title: Moses, venerable
\latin{in typo sanctae Ecclesiae}, gives thanks that God was his refuge from his
exposure as an infant to the murmuring of the people, and he speaks for every
church, since to the whole Church, \latin{quae in hac saeculi tempestate
versatur}, the Lord is always the safest refuge (\work{Expositio Psalmorum}
89.1). The exposition of the psalms printed under the name of St Athanasius
finds in the first verse a proem drawn from benefits already given, under
Pharaoh, under Josue and in Babylon, and in the second the confession that the
Lord coexisted with the Father before all creation, \latin{eum ante omnia
opificia Patri coexsistere} (PG~27, col.~395). Durandus says that the soldier
wearied or wounded has his refuge in God alone (VI.138). Fromage hears in the
Gradual and the Alleluia the Church preparing its last exodus from the world
(p.~463), a phrase the third reading takes up; Schuster, the assurance that the God who helped Israel and
the early Church is ``still the same God, faithful and unchanging'' (p.~180).

\subsection{Alleluia: Israel out of Egypt}

The verse is the first of the psalm \latin{In exitu Israel}, which remembers the
Exodus and the entry into the land: ``The sea saw and fled: Jordan was turned
back. The mountains skipped like rams'' (Ps 113:3--4). In the Vulgate the psalm
runs on into what the Hebrew numbers as Ps~115. The verse names the two ends of
the going out: Israel, the house of Jacob, and Egypt, ``a barbarous people''.
Bellarmine explains that on leaving Egypt the Jews began to be a people under
God's special rule, and that the Egyptians are called barbarous because they
spoke a strange language. Schuster notes that Egypt was far richer in civil
culture than Israel, and that Scripture calls it barbarous by the measure of the
spirit; he adds that the psalm was among those sung at the Passover meal
(p.~180).

St Augustine teaches that the deeds of the Exodus ``happened unto them for
ensamples'' (1~Cor 10:11). Egypt, a name he understands as affliction or
oppression, stands for this world, from which no one withdraws in heart without
God's mercy, as Israel did not leave Egypt without his help (\work{Enarrationes}
113.1--4; the third reading gives his words). Cassiodorus takes the going out as the moment when we are released
from the bonds of sin, \latin{a peccatorum vinculis}, and freed from the host of
the Egyptians, that is, of the demons; we truly become Israel when the pomps of
this world begin to leave our minds (\work{Expositio Psalmorum} 113.1). St John
Chrysostom, expounding the same verse, observes that God first shows his power
and his benefits, freeing his people from a cruel enemy, and only then asks for
worship, where men wish first to rule and only afterwards to do good
(\work{Expositio in Psalmum} 113.1).

\subsection{Gospel: the king's reckoning}

The parable closes the discourse of Mt~18, spoken after Jesus and the disciples
have come to Capharnaum (17:23) and before he leaves Galilee (19:1). The
chapter begins with the question who is greatest in the kingdom and goes on to
the little ones, scandal, the lost sheep and the brother who sins, who is to be
rebuked first alone, then before witnesses, then before the Church. Peter then
asks how often he must forgive his brother, ``till seven times?'', and hears
``till seventy times seven times'' (18:21--22). ``Therefore'', the Vulgate
goes on, ``is the kingdom of heaven likened to a king''. The Mass begins
the parable without Peter's question, so that it stands on its own as a word of
Jesus to his disciples.

The story moves in three scenes. In the first the king begins his reckoning,
\latin{voluit rationem ponere}, and a servant is brought in who owes ten thousand
talents, a sum no servant could repay; the second reading gives the hundred
pence's value and Augustine's reckoning of the difference. The king orders him sold with his wife,
his children and all he has; the servant falls down and asks, ``Have patience
with me, and I will pay thee all''; and the lord, moved with pity, lets him go and
forgives the debt. He asked for time and was given remission. In the second the
same servant goes out and finds a fellow servant who owes him a hundred pence,
takes him by the throat, \latin{suffocabat}, and hears the very words he had
spoken, ``Have patience with me, and I will pay thee all''; and he will not, but
casts him into prison until he pays. In the third the fellow servants, grieved,
tell their lord; the lord calls the servant ``wicked'', reminds him that he was
forgiven ``because thou besoughtest me'', asks whether he should not have had
compassion as he had received it, and in anger delivers him ``to the torturers
until he paid all the debt''. The sentence that closes the Gospel applies the
whole to the hearers: ``So also shall my heavenly Father do to you, if you
forgive not every one his brother from your hearts.''

St Jerome explains that it is \latin{familiare \dots\ Syris et maxime
Palaestinis} to join a parable to what they say, so that a precept the hearer
cannot hold by itself is held by its likeness; and he states the parable's
argument as a question: \latin{Si \dots\ dominus servo debitori decem millium
talentorum tam facile dimisit, quanto magis servi conservis suis debent minora
dimittere?}, if the lord so easily forgave his servant ten thousand talents, how
much more ought servants to forgive their fellow servants lesser debts?
(\work{Commentarii in Matthaeum} III, PL~26, cols.~132--133). St Hilary of
Poitiers says that the parable teaches pardon without measure through the grace
of the Gospel, and that the Lord himself expounded its whole meaning
(\work{Commentarius in Matthaeum} 18.10--11). The liturgical commentators who
reach this Gospel read it toward the judgment and the Lord's Prayer: Rupert hears
in the servant's debt the petition \latin{Dimitte nobis debita nostra}
(\work{De divinis officiis} XII.22); Durandus says that on the day of judgment the
Lord will settle accounts with us (\work{Rationale} VI.139); Honorius, that
Christ the true king will come then to judgment and require of each his deeds
(\work{Gemma animae} IV.93). Schuster calls the parable ``the most authoritative
commentary on those words of the Lord's Prayer: `Forgive us our trespasses as we
forgive them that trespass against us''' (p.~180).

\subsection{Offertory: the man of Hus}

The Offertory tells the first two chapters of Job in a narrative of its own. In
the book, Job of the land of Hus, which the Lamentations join to Edom (Lam
4:21), is ``simple and upright, and fearing God, and avoiding evil'' (Job 1:1).
When the sons of God come to stand before the Lord, Satan comes among them; the
Lord asks, ``Hast thou considered my servant, Job?'', and Satan answers, ``Doth
Job fear God in vain?'' The Lord gives him leave over all that Job has, ``only
put not forth thy hand upon his person'' (1:6--12). Job loses his oxen, his
asses, his sheep, his camels, his servants and his sons and daughters, and says,
``the Lord gave, and the Lord hath taken away \dots\ blessed be the name of the
Lord. In all these things Job sinned not by his lips'' (1:21--22). At a second
council Satan is given Job's flesh, ``but yet save his life'' (2:6), and strikes
him ``with a very grievous ulcer, from the sole of the foot even to the top of
his head'' (2:7). The antiphon gathers this into a single movement: the just man,
the adversary's petition, the power given by the Lord, the loss of goods and
sons, the wound in the flesh. It says that Satan \latin{petiit}, asked, a verb
the book does not use, and it says nothing of Job's answer.

St Thomas Aquinas, expounding the book according to its literal sense, holds
that it tells a thing that happened and not a parable; that Satan can harm
\latin{non quantum vult sed quantum permittitur}, not as much as he wills but as
much as he is permitted; and that the will to harm is his own while the power
comes only from God (\work{Expositio super Iob} 1). St Gregory the Great, whose
\work{Moralia} on Job is the great Latin commentary, finds two wills in the
trial: the devil asked for it that Job might fall short, the Lord in kindness
permitted it, and Job went forward instead (Preface~7). The first reading follows that answer.

Fromage makes the Offertory the centre of the Sunday: it ``contains the ruling
idea, which runs through this twenty-first Sunday after Pentecost'' (p.~455). The
other commentators expound the antiphon itself. Durandus says that the devil is
conquered above all by patience (VI.138); Rupert, that what befell Job in history
befalls the soldier of the virtues in his moral life, when the tempter, permitted
by a hidden dispensation of God, strikes him (XII.21); Honorius, that Job
signifies Christ, who suffered for us (IV.90--91); Sicard of Cremona, that Job is
the example of the brave warrior (\work{Mitrale} VIII.21); for Schuster, Job ``is
considered as being a figure of Christ crucified'' (p.~178). Fromage and Schuster both print or describe the
verses of Job's lament that the older chant books sang after the antiphon; the
1962 and 1862 Missals print the antiphon alone.

\subsection{Secret: the sacrifice by which God willed to be appeased}

The Secret asks God to receive the offerings \latin{propitius}, graciously, and
names what they are: the sacrifices by which he willed to be appeased,
\latin{quibus et te placari voluisti}, and by which he willed salvation to be
restored to us by his powerful mercy. The prayer makes the appeasing God's own
will, and the restoring the work of his \latin{pietas}. Schuster explains that
the Eucharistic sacrifice both satisfies divine justice for us and restores us
to the salvation in which God created our first parents, and that the attribute
\latin{potens} given to God's mercy fits a mercy that ``has triumphed over divine
justice itself'' (Jas 2:13; p.~181). Fromage says that the salvation of the
world and of each man ``is, virtually, ever in the august Sacrifice, whose power
restores man, by appeasing God'' (p.~470).

\subsection{Communion: the plea of the strophe Caph}

The eleventh strophe of Ps~118, each of whose verses begins in Hebrew with the
letter \latin{Caph}, is the lament of a man worn out with waiting. His soul
faints for God's salvation; his eyes fail for God's word, ``When wilt thou
comfort me?''; he has become ``like a bottle in the frost''; he asks, ``How many
are the days of thy servant: when wilt thou execute judgment on them that
persecute me?''; the wicked have told him fables, and have persecuted him
unjustly; ``They had almost made an end of me upon earth: but I have not forsaken
thy commandments''; ``Quicken thou me according to thy mercy'' (vv.~81--88).
The Communion keeps from this strophe the hope in God's salvation and word
(v.~81, without its verb \latin{Defecit}), the question when God will judge the
persecutors (v.~84, without its first question), and the cry for help against
the wicked (v.~86, without its first half), and adds \latin{Domine Deus meus}.
What it sings is hope, a question about the judgment, and a plea for help; the
fainting, the counting of days and the near death of the strophe are left
unsung.

Who the persecutors are, and when the judgment asked for will come, are the
questions on which the Fathers read the verse most differently, and the
readings below take them up at length. The liturgical commentators hear in the
Communion, first, hope: Durandus says that patience avails nothing without hope
(VI.138); Fromage, that ``An unflagging hope ever accompanies the admirable
patience of holy Church'', so that persecutions ``never interrupt her prayer''
(p.~470). Honorius, reading it in the history of Esther, hears thanks that God
has done judgment on the persecutors (IV.90). Schuster says that the spiritual
enemies whose judgment is asked are judged already, by their own refusal of the
light, ``Qui non credit, jam judicatus est'' (Jn 3:18; p.~181).

\subsection{Postcommunion: the food of immortality}

The Postcommunion calls what has been received \latin{immortalitatis alimonia},
the food of immortality, and asks that what the mouth has received the pure
mind may follow, \latin{pura mente sectemur}. Its pairing of the mouth and the
mind, of the outward reception and the inward pursuit, is the pairing of the
Order of Mass's \latin{Quod ore sumpsimus}. Schuster calls the Eucharist ``the
food which nourishes us for all eternity'' and reads the petition as asking that
sacramental Communion bind ``our mind, our heart, and our will to Jesus''
(p.~181); Fromage: ``Now that we have been nourished by the food of immortality,
let us live on it, with all the evidence of a soul that is made pure''
(p.~470).
